\section{Proof of the Bound \texorpdfstring{$\Delta^2+\lfloor\frac{\Delta}{2}\rfloor^2$}{Delta squared plus floor(Delta/2) squared}}\label{sec:exact}\label{sec:multi-odd}

Put $D=\Delta(G)$. If $D$ is even, then $f(D)=\frac{5}{4}D^2$, so
\textbf{Lemma~\ref{prop:multi-uniform}} gives the asserted bound,
including $D=0$. To set up the proof for odd $D\geq3$, suppose that
$G$ has a strong clique with more than $f(D)$ edges. Mark
exactly $m=f(D)+1$ of these edges and restrict $G$ to their endpoints.
For two disjoint marked edges, both endpoints of any edge joining them
are retained. Hence the marked edges still form a strong clique.
Retain all marked edges. For each pair of adjacent vertices joined by
no marked edge, retain one unmarked edge. Discard the other unmarked
edges. The resulting multigraph has degree bound $D$ and induces the
same adjacencies on the marked endpoints.

As in \textbf{Subsection~\ref{sec:weighted-counts}}, let $s$ be the
number of distinct pairs of endpoints of marked edges and put $b=m-s$. Write
$S=\frac{s}{D^2}$, $M=\frac{m}{D^2}$, and $e=\frac{b}{m}$, so
$S=M(1-e)$. Since $f(D)=\frac{5D^2-2D+1}{4}$ for odd $D$, our
choice of $m$ gives
\begin{equation}\label{eq:multi-mass}
 M=\frac{5}{4}-\frac{1}{2D}+\frac{5}{4D^2}.
\end{equation}
Let $G_{\rm s}$ denote the simple support. Color an edge of $G_{\rm s}$
red if its endpoints are joined by a marked edge.

\subsection{Excluding Counterexamples for \texorpdfstring{$\Delta\geq9$}{Delta >= 9}}\label{subsec:exact-large-degree}

Assume that $D\geq9$ is odd. We first show that a counterexample
has two parallel marked edges and that $G_{\rm s}$ contains an
induced five-cycle with all five edges red. Degree counts around
these cycles then restrict the possible values of $e$. On the
resulting interval, the weighted flag inequalities contradict
\eqref{eq:multi-mass}.

For each induced red five-cycle $c_0c_1c_2c_3c_4c_0$ in $G_{\rm s}$,
take indices modulo five. For each $i$, let $U_i$
consist of the vertices whose neighbors on the cycle are exactly
$c_{i-1},c_{i+1}$, with both adjacencies red. In particular $c_i\in U_i$.
If $x\in U_i$ and $y\in U_{i+1}$, then the disjoint red edges
$xc_{i-1}$ and $yc_{i+2}$ must have a joining edge. The definitions
of the two sets exclude $xc_{i+2}$ and $yc_{i-1}$, and the cycle
has no chord $c_{i-1}c_{i+2}$. Hence $xy\in E(G_{\rm s})$.
Set
\[
 B=\bigcup_{i=0}^4 U_i,\qquad
 s_i=D-|U_{i-1}|-|U_{i+1}|,\qquad
 R=\sum_{i=0}^4s_i=5D-2|B|.
\]
At $c_i$ the adjacencies to $U_{i-1}\cup U_{i+1}$ use at least
$|U_{i-1}|+|U_{i+1}|$ of its degree, so $s_i\geq0$. Since $D$ is
odd, $R=5D-2|B|$ is odd.

Let $q_C$ be the normalized count of ordered induced five-cycles
with all five edges red in the simple support, using $2sD^3$ as the
denominator. Let $L_R$ be the sum of $\frac{R}{D}$ over these ordered cycles,
divided by the same denominator. The simple support can be completed
to a $D$-regular simple graph as in~\cite[Section~G.1]{FGSTSupplement}.
The old vertices
induce their original graph, and no added edge is red. Thus the
cycles under consideration and the sets $U_i$ of each such cycle
are preserved. The simple graph inequalities for $S,q_C,L_R$
in~\cite[(G.1)]{FGSTSupplement} apply. If $R\geq\alpha D$ on
every induced red five-cycle, then $L_R\geq\alpha q_C$ by summing
over the cycles.

\begin{lemma}\label{lem:extra-copy-required}\label{lem:counterexample-cycle}
Suppose that $D\geq9$ is odd and that $m=f(D)+1$ marked edges form
a strong clique. Then $b>0$ and $q_C>0$.
\end{lemma}

\begin{proof}
For $D\geq9$, the value of $M$ in \eqref{eq:multi-mass} satisfies
\[
 M-\frac{98}{81}=\frac{(D-9)(13D-45)}{324D^2}\geq0.
\]
If $q_C=0$, the bound in~\cite[(G.2)]{FGSTSupplement} gives
$M<\frac{98}{81}$. Thus $q_C>0$.

At $D=9$, suppose $b=0$. The $m$ marked edges then lie
on $m$ distinct pairs in the simple support. They form a strong
clique there, so \textbf{Lemma~D.1} of the supplement~\cite{FGSTSupplement}
gives $m\leq f(D)$, contrary to $m=f(D)+1$.

For odd $D\geq11$, the degree estimates
in~\cite[Section~E.4]{FGSTSupplement} give $R>\frac{D+1}{2}$ on every
induced red five-cycle. These estimates count marked incidences in
the reduced multigraph and apply for every $b\geq0$. In particular,
$L_R\geq\frac{q_C}{2}$. The last inequality
in~\cite[(G.1)]{FGSTSupplement} now gives
\[
 S\leq\frac{5}{4}-\frac{0.45119}{16}+2\cdot10^{-9}
   =1.221800627.
\]
For $D\geq15$ we have $M\geq\frac{11}{9}$, since
$M-\frac{11}{9}=\frac{(D-15)(D-3)}{36D^2}$ by
\eqref{eq:multi-mass}. Consequently $S<M$ and $b>0$.

For $D=11$, the degree count in~\cite[Section~E.5]{FGSTSupplement}
also excludes $R=7$, so every induced red five-cycle has $R\geq9$.
For $D=13$, the strict inequality $R>7$ and the odd parity of $R$
give the same lower bound. Thus $L_R\geq\frac{9}{D}q_C$ in both
cases. Put $a=0.99999$, $c=0.5682\cdot\frac{9}{D}$ and
$\varepsilon=10^{-9}$. The first two inequalities
in~\cite[(G.1)]{FGSTSupplement} give
\[
 S\leq\frac{5}{4}-\frac{a}{8}+2\varepsilon+a q_C,
 \qquad
 S\leq\frac{5}{4}+\varepsilon-c q_C .
\]
Multiply the first inequality by $c$ and the second by $a$.
Adding cancels $q_C$ and gives
\[
 S\leq\frac{5}{4}-\frac{ac}{8(a+c)}
             +\frac{(2c+a)\varepsilon}{a+c}.
\]
At $D=11$ this upper bound is less than $\frac{1211}{1000}$, whereas
$M=\frac{147}{121}$. At $D=13$ it is less than $\frac{1215}{1000}$,
whereas $M=\frac{206}{169}$. In both cases $S<M$, as required.
\end{proof}

By \textbf{Lemma~\ref{lem:counterexample-cycle}}, we may choose two
parallel marked edges, and the simple support contains an induced red
five-cycle. For every such cycle the calculation
in~\cite[Section~E.5]{FGSTSupplement} gives $R\geq9$ if
$D=11$ and $R\geq11$ if $D\in\{13,15\}$. The degree count at the
endpoints of two parallel marked edges
in~\cite[Section~F.2]{FGSTSupplement} gives $R\geq11$ at $D=9$.
For odd $D\geq11$, the inequality $R>\frac{D+1}{2}$ together with the
odd parity of $R$ gives
\begin{equation}\label{eq:multi-slack}
 R\geq
 \begin{cases}
 \frac{D+5}{2} & \text{if }D\equiv1\pmod4,\\
 \frac{D+3}{2} & \text{if }D\equiv3\pmod4.
 \end{cases}
\end{equation}

Let $[M_-,M_+]$ and $[e_-,e_+]$ be the ranges for $M$ and $e$
in \textbf{Table~\ref{tab:multi-odd}}, respectively. The fourth
column is the lower bound $\alpha$ on $\frac{R}{D}$ for every induced
red five-cycle. In the last column we round the verified upper bounds
upwards to short rationals. The last certificate gives
infeasibility directly. The exact fractions and row counts are in the
companion record \texttt{parity\_family\_replayed.json}.

\begin{table}[t]
\centering
\small
\begin{tabular}{@{}lcccc@{}}
\toprule
Odd $D$ & $[M_-,M_+]$ & $[e_-,e_+]$ & $\alpha$ & Verified conclusion\\
\midrule
$9$        & $[1.209,1.210]$  & $[0.020,0.180]$   & $1.222$ & $M<-1.973$\\
$11$       & $[1.2148,1.215]$ & $[0.0135,0.360]$  & $0.818$ & $M<0.785$\\
$13$       & $[1.2185,1.219]$ & $[0.009,0.362]$   & $0.846$ & $M<0.573$\\
$15$       & $[1.222,1.223]$  & $[0.007,0.365]$   & $0.733$ & $M<0.827$\\
$17$       & $[1.2249,1.225]$ & $[0.006,0.368]$   & $0.647$ & $M<0.631$\\
$19$       & $[1.227,1.228]$  & $[0.006,0.370]$   & $0.578$ & $M<0.910$\\
$21$--$23$ & $[1.2249,1.231]$ & $[0.004,0.372]$   & $0.565$ & $M<1.227683$\\
$25$--$37$ & $[1.232,1.238]$  & $[0.009,0.384]$   & $0.542$ & $M<0.177$\\
$39$--$47$ & $[1.238,1.240]$  & $[0.0134,0.382]$  & $0.531$ & $M<-0.676$\\
$49$--$57$ & $[1.240,1.242]$  & $[0.0082,0.380]$  & $\frac{1}{3}$ & $M<0.292$\\
$\ge59$    & $[1.24175,1.25]$ & $[0.005,0.250]$   & $0.228$ & infeasible\\
\bottomrule
\end{tabular}
\caption{Parameter intervals and verified bounds for odd $D$. All decimal entries denote rational numbers.}
\label{tab:multi-odd}
\end{table}

The quantity in \eqref{eq:multi-mass} is increasing for $D\ge5$ and is
less than $\frac{5}{4}$. Evaluating it at the endpoints of each degree
range gives the corresponding bounds on $M$.
For $D=9$ we have $2\le b\le17$ by~\cite[Section~F.2]{FGSTSupplement}, so
$e\in[\frac{2}{98},\frac{17}{98}]$ and $S\le\frac{32}{27}$. For $D=11,13$,
the reduction in~\cite[Section~E.5]{FGSTSupplement} excludes $b=1$. Thus $S$ is at most
$\frac{145}{121}$ and $\frac{204}{169}$, respectively, giving the required lower
bounds on $e$. In the remaining rows the simple projection bound
$S\le S_{\max}$ obtained from $\alpha$
in~\cite[Section~G.1]{FGSTSupplement} gives $e\ge1-\frac{S_{\max}}{M}$. For an
interval ending at $D_+$, recall that $j_v$ counts the extra marked
copies incident with $v$ and put $A=\frac{\sum_v j_v^2}{2bD}$.
\textbf{Lemma~\ref{lem:multi-extra}} gives
\[
 A\le A_+:=\frac{2-M_--D_+^{-2}}{2+D_+^{-1}}.
\]
We use $e\le1-\frac{M_-}{2}$ or
$e\le\frac{2A_+(1-A_+)-\frac{A_+}{D_+}}{M_-}$ when the latter is smaller.
The function $A\mapsto 2A(1-A)-\frac{A}{D_+}$ is increasing for
$0\le A\le A_+$ in these intervals.
For $D\ge59$ we use \textbf{Lemma~\ref{prop:multi-uniform}} to bound $e$.
The rational replay checks these interval comparisons.

The finite-degree constraints for $D\le23$ also use the exclusion of
a red $K_4$ proved in~\cite[Section~F.3]{FGSTSupplement} and, for specific $D$,
the rooted triangle and repeated-pair restrictions
in~\cite[Section~G]{FGSTSupplement}.
For each row in Table~\ref{tab:multi-odd}, the weighted graph counts must
satisfy the degree equations, the positive semidefinite constraints,
and the inequalities for that row. The verified rational linear
combination of these constraints gives the conclusion in the last
column, as detailed in~\cite[Section~G.4]{FGSTSupplement}. For the row $21$--$23$, for instance, the
upper bound $1.227683$ is strictly smaller than $\frac{542}{441}$, the lower
bound from \eqref{eq:multi-mass}. On every other finite degree
interval in the table, the verified upper bound also lies below
the least value of $M$ given by \eqref{eq:multi-mass}. For
$D\geq59$, the last certificate rules out every set of counts
satisfying the constraints.

\subsection{The Cases \texorpdfstring{$\Delta<9$}{Delta < 9}}\label{subsec:exact-small-degree}

The remaining odd values of $D$ are $1,3,5,7$. If $D=1$, each
component has at most one edge, so $\omega(L(G)^2)\leq1=f(1)$.
Suppose that $D\in\{3,5,7\}$ and that $m=f(D)+1$ marked edges form
a strong clique. If $b=0$, these edges give $m$ distinct edges in
the simple support $G_{\rm s}$. They still form a strong clique
because any two disjoint marked edges have endpoints joined by an
edge of $G$, and those endpoints are also adjacent in $G_{\rm s}$. Since
$\Delta(G_{\rm s})\leq D$, \textbf{Lemma~D.1} of the supplement~\cite{FGSTSupplement}
gives $m\leq f(D)$, a contradiction. Thus two of the marked edges
are parallel. The reductions in~\cite[Section~F.1]{FGSTSupplement} exclude
a strong clique of size $f(D)+1$ containing such a pair for each of
$D=3,5,7$.

\begin{lateproof}{\textbf{Theorem~\ref{thm:multigraph}}}
For even $D$, the result follows from \textbf{Lemma~\ref{prop:multi-uniform}}.
\textbf{Subsection~\ref{subsec:exact-large-degree}} excludes a
counterexample for every odd $D\geq9$, and
\textbf{Subsection~\ref{subsec:exact-small-degree}} proves the bound
for each smaller odd $D$.
\end{lateproof}
